\section{Discussion and limitations}

The results provide a population benchmark for passive outcome annotation from a common distribution. Under the independent sampling condition in \cref{obj:ass:data-product}, complete records supply outcomes together with treatment and covariates, while auxiliary records supply treatment and covariates. The minimax comparison in \cref{obj:thm:uniform-annotation-frontier} characterizes the precision of the population contrast in \cref{obj:def:target} as both budgets, the covariate alphabet, and the public overlap floor vary. Under \cref{obj:ass:consistency,obj:ass:exchangeability}, this contrast is the population average treatment effect.

The marginal-budget plateau gives these information sources a distinct economic interpretation. Additional auxiliary records improve knowledge of treatment--covariate frequencies and reduce the many-cell component of the risk. Once that component is at or below the rare-label benchmark, the minimax risk has order \(\min\{1,(n\epsilon)^{-1}\}\), by \cref{obj:thm:uniform-annotation-frontier,obj:thm:rare-label-floor}. Thus auxiliary expansion can bring estimation to a precision scale governed by rare-arm outcomes in the complete sample. On the oracle sequences characterized in \cref{obj:thm:annotation-resource-phases}, this scale is eventually \((n_k\epsilon_k)^{-1}\). The plateau concerns risk order: it identifies the information budget governing precision once treatment--covariate information is sufficiently abundant.

The exact-marginal experiment sharpens this interpretation. Supplying the entire treatment--covariate probability table yields rare-label risk order for every allowed alphabet size, as established in \cref{obj:prop:benchmark-recovery}. Moreover, for each fixed complete-record sample size \(n\ge1\), fixed alphabet size \(d\ge2\), and fixed overlap floor \(0<\epsilon\le1/4\), the minimax risk converges to the exact-marginal risk as the integer auxiliary budget tends to infinity. This fixed-parameter limit describes what increasingly precise frequency information can deliver with a given outcome sample. The common-marginal alternatives in \cref{obj:thm:rare-label-floor} explain the remaining statistical difficulty: identical treatment--covariate frequencies accompany different rare-arm outcome distributions and different target values.

For annotation planning, these comparisons distinguish investment in population frequency information from investment in outcome information. They provide a benchmark for the precision attainable when acquisition follows the common-distribution experiment. Allocation procedures that choose which records receive outcomes introduce a further design problem. The asymptotic-variance allocation and feasible batch-adaptive procedure of \citet[Theorems 4.2 and 5.1 of the arXiv manuscript]{NwankwoGoldkindZhou2026Annotations} provide complementary context for that problem under their stated budget, identification, positivity, and regularity conditions.

\subsection{Limitations and future work}

The attaining estimator in \cref{obj:def:estimator-handle} outputs zero whenever \(n\epsilon<\exp(4096)\). Its uniform result therefore certifies minimax order with proof-oriented constants rather than providing calibrated finite-budget tuning: the comparison constant absorbs this threshold over all realistic values below it. The resource conditions in \cref{obj:thm:annotation-resource-phases} should accordingly be read as order-based planning benchmarks, not as calibrated recommendations for a particular annotation budget. Computationally practical tuning, stable evaluation, and implementation at moderate sample sizes require further analysis. Adaptive allocation would also change the observation mechanism in \cref{obj:ass:data-product}; extending the benchmark to allocation rules would require accounting for the dependence and selection induced by those rules.

The paper's loss is squared error for a population treatment contrast. Interval estimation requires additional distributional and coverage guarantees, while policy objectives require specifying the decision class and the welfare criterion. These questions have different statistical requirements from the minimax point-estimation comparison.

Finally, the overlap floor is a public restriction on population propensities in occupied cells. Empirical overlap assessments introduce uncertainty about that restriction, particularly when observed cell counts are small. Extensions to continuous covariates, multiple treatments, or outcomes beyond the binary case would require corresponding model and information-budget characterizations. The binary finite-alphabet experiment provides a reference point for studying these extensions.
